\chapter{Facial Droop}

\section*{Definition}
Unilateral or bilateral weakness of facial muscles resulting in asymmetry of expression, inability to close the eye or raise the eyebrow, and flattening of the nasolabial fold.

\section*{What Happens in Your Body}
Facial droop results from dysfunction of the facial nerve (cranial nerve VII) at any point from its nucleus in the pons to its peripheral branches. Lower motor neuron lesions (Bell palsy, Ramsay Hunt, Lyme) affect entire ipsilateral face; upper motor neuron lesions (stroke, tumor) spare the forehead due to bilateral cortical innervation.

\section*{Causes}
\begin{itemize}
  \item Bell palsy (idiopathic facial nerve palsy)
  \item Stroke or transient ischemic attack (upper motor neuron)
  \item Ramsay Hunt syndrome (herpes zoster oticus)
  \item Lyme disease (borreliosis)
  \item Brain tumor or metastatic lesion
  \item Traumatic facial nerve injury
  \item Otitis media or mastoiditis
  \item Guillain-Barre syndrome (bilateral facial weakness)
  \item Sarcoidosis (Heerfordt syndrome)
  \item Myasthenia gravis (fatigable weakness)
\end{itemize}

\section*{When to See a Doctor}
See your doctor if:
\begin{itemize}
  \item Symptoms are severe or getting worse over time
  \item Sudden facial droop requires immediate medical evaluation to differentiate stroke from Bell palsy (FAST assessment)
  \item You have other concerning symptoms such as fever, weight loss, or bleeding
\end{itemize}

\section*{Related Symptoms}
\begin{itemize}
  \item Arm or leg weakness
  \item Speech difficulty
  \item Ear pain
  \item Eye closure difficulty
  \item Drooling
\end{itemize}
\textbf{Important:} If this symptom persists, your doctor will check for serious underlying causes.

\section*{Self-Care / Home Management}
\begin{itemize}
  \item Protect the eye on the affected side with lubricating drops and a patch, especially during sleep.
  \item Chew on the unaffected side and eat soft foods if chewing is difficult.
  \item Perform gentle facial exercises in front of a mirror to maintain muscle tone.
  \item Avoid exposure to cold wind and wear a scarf over the face when outdoors.
  \item Seek immediate emergency evaluation for sudden facial droop to rule out stroke.
\end{itemize}

\section*{How Your Doctor Will Evaluate You}
\begin{itemize}
  \item Your doctor will perform a FAST assessment (Face, Arms, Speech, Time) to screen for stroke.
  \item A detailed neurological examination determines whether the lesion is upper or lower motor neuron.
  \item CT or MRI brain imaging is performed emergently if stroke is suspected.
  \item Blood tests screen for infectious causes (Lyme titers, HSV/VZV serology).
  \item Electromyography (EMG) may be ordered to assess facial nerve function and prognosis.
\end{itemize}

\section*{Treatment Options}
\begin{itemize}
  \item Oral corticosteroids within 72 hours for Bell palsy to improve recovery.
  \item Antiviral therapy (acyclovir, valacyclovir) if Ramsay Hunt syndrome is suspected.
  \item Antiplatelet or anticoagulant therapy for stroke, directed by stroke severity and timing.
  \item Eye protection (lubricating drops, taping the eye closed at night) for lagophthalmos.
  \item Facial physiotherapy and biofeedback for persistent facial weakness.
\end{itemize}


\section*{References}
\begin{thebibliography}{99}
\bibitem{jauch2008}
Jauch EC, Saver JL, Adams HP Jr, et al. Guidelines for the early management of patients with acute ischemic stroke. \textit{Stroke}. 2013;44(3):870--947.
\bibitem{holland2008}
Holland NJ, Weiner GM. Recent developments in Bell palsy. \textit{BMJ}. 2004;329(7465):553--557.
\bibitem{sullivan2007}
Sullivan FM, Swan IRC, Donnan PT, et al. Early treatment with prednisolone or acyclovir in Bell palsy. \textit{N Engl J Med}. 2007;357(16):1598--1607.
\bibitem{gilden2001}
Gilden DH. Bell palsy. \textit{N Engl J Med}. 2004;351(13):1323--1331.
\bibitem{bersani2017}
Bersani G, Iannitelli A, Pacitti F, et al. Facial nerve palsy: a review. \textit{Clin Ter}. 2017;168(3):e182--e189.
\bibitem{engstrom2008}
Engstr\"{o}m M, Berg T, Stjernquist-Desatnik A, et al. Prednisolone and valacyclovir in Bell palsy: a randomised, double-blind, placebo-controlled study. \textit{Lancet Neurol}. 2008;7(11):993--1000.
\end{thebibliography}
